\subsection{The matching engine pair}
\label{sec:ha:enabled:me}

Two matching engines run. One acts; the other holds a replica of the set of open
orders, built from the same records in the same order.

\subsubsection{What state it holds}
\label{sec:ha:enabled:me:state}

The set of open orders and the sequence number it is current to ---
\cref{sec:ha:disabled:me:state} --- held twice. The engine keeps no log of its
own in either configuration: the record of what was taken belongs to the
sequencer, and what the engine holds is the consequence of it.

\subsubsection{Where the copy that outlives it lives}
\label{sec:ha:enabled:me:durable}

In the other instance, kept current by the same stream. This is the copy that
\cref{sec:ha:disabled:me:durable} says does not exist when the configuration is
disabled, and its absence there is why a checkpoint is needed at all.

\subsubsection{How a restarted instance becomes current}
\label{sec:ha:enabled:me:recovery}

It reports the sequence number its replica has reached, receives everything
committed after it, applies that silently, and begins acting.

\textbf{This is the same operation as \cref{sec:ha:disabled:me:recovery}}, and
the subsection is short for that reason. A restarting instance learns its
starting number from a checkpoint; a promoted instance knows it from the replica
it has been maintaining. Nothing else about the operation differs, which is why
the disabled configuration is worth getting right first: it builds the mechanism
that this one uses.

The engine differs from the sequencer in one respect. A sequencer pair settles
leadership between themselves in both directions, because both hold both
generations at the same moment. A matching engine only ever \emph{gives}
leadership up: it adopts the follower position on its peer's word and takes the
generation it is given, and never issues one. Deferring cannot create a second
instance acting, whatever the peer is confused about, which is what makes acting
on the peer's word alone safe.

\subsubsection{What the member is told}
\label{sec:ha:enabled:me:member}

\begin{requirement}{R-0073}{An open order survives a failover}
  An order open before leadership moved shall be open after it, on the same
  terms, and shall remain cancellable by the member that placed it. Where the set
  cannot be recovered, \reqref{R-0020} applies instead and every order the venue
  can still name is cancelled to its member.
  \rationale{\reqref{R-0018} requires this of a restart. A member that keeps its
  orders when a process dies and loses them when a machine dies is being given a
  worse outcome by the configuration that exists to give it a better one.

  The two requirements are a hierarchy rather than alternatives: recover what can
  be recovered, and disclaim what cannot. The member is still told exactly
  where it stands, which was the merit of cancelling, but only in the case where
  the venue does not have the orders.}
  \uncovered{scenario 21 asserts the fallback, and nothing asserts this}
\end{requirement}

\begin{gap}{}
  A promoted matching engine cancels the orders that were open, whether or not it
  could have recovered them. Under the requirement above that is the fallback
  path, and the venue takes it every time. The behaviour was settled before
  promotion was known to recover the set reliably --- which it has since been
  measured to do --- and was not revisited afterwards.
\end{gap}

\begin{requirement}{R-0074}{An instance that gives up leadership cannot thereby create a second one}
  An instance shall adopt the follower position only on its peer's assertion that
  the peer leads, only while itself holding no position, and only where the
  generation offered is not older than one it has already seen.
  \rationale{Giving up a position cannot produce two instances acting, whatever
  the peer is mistaken about, which is what makes a one-way announcement
  sufficient here where the sequencer needs an exchange.}
  \uncovered{no scenario offers a stale generation to an instance holding no position}
\end{requirement}

\begin{requirement}{R-0075}{The venue keeps working when the acting engine is lost}
  Members shall continue to be able to place and cancel orders across the loss
  of the matching engine instance that was acting.
  \rationale{Cancelling matters here as much as placing, and is the half that
  is easier to overlook: a member that can place orders but cannot cancel the
  ones it already has is worse off than one that can do neither, because it is
  accumulating exposure it has no way to reduce.}
  \covers{ha_test:16}
\end{requirement}

\begin{requirement}{R-0076}{Orders taken while no engine was acting are applied on promotion}
  An order committed while no instance was acting shall be applied and answered
  once one is.
  \rationale{The venue took those orders, so it must answer each of them. This is
  the case the disabled configuration cannot serve, because there is no instance
  holding a position from which to catch up.}
  \uncovered{no scenario asserts that deferred orders are answered after a promotion}
\end{requirement}
